% GENERATED FILE. Edit canonical lessons, then run npm run generate:books.
% canonical-source-hash: fnv1a64:7685f1759d3057bf
% canonical-lessons: MR-C204-gambhir, MR-C204-lajalu, MR-C204-sur, MR-C204-bhitra, MR-C204-utsuk

\chapter{Serious, Shy, Valiant, Fearful, Curious}
\label{ch:mr-serious-shy-valiant-fearful-curious}
\hlchaptermodality{204}

\begin{chapteropening}
\textbf{By the end of this chapter:} \emph{I can say serious, shy, valiant, fearful and curious.}

Complete the last lesson of chapter 204: I can say serious, shy, valiant, fearful and curious.
\end{chapteropening}

\section[\texorpdfstring{\mr{गंभीर} (gambhīr)}{गंभीर (gambhīr)}]{\mr{गंभीर} (gambhīr) — serious}

\label{lesson:MR-C204-gambhir}



\begin{quote}
Before the new one: say the Marathi for about, then the Marathi for only.
\end{quote}

\subsection*{You'll want to know: \mr{गंभीर}}
\textbf{\mr{गंभीर}} — \emph{gambhīr} — \textquotedblleft{}serious\textquotedblright{}.

\begin{cousinweb}
\textbf{\mr{गंभीर}} (\emph{gambhīr}) is Sanskrit \emph{gambhīra}, \textquotedblleft{}deep\textquotedblright{}, taken over whole as a learned loan, a tatsama (CDIAL 4031). CDIAL lists the Marathi word as \emph{gãbhīr}, \textquotedblleft{}deep (of water), serious\textquotedblright{}.
\end{cousinweb}

\begin{grammarlens}[title={What you've built}]
One more word for this chapter.
\end{grammarlens}

\subsection*{Your turn}
\begin{itemize}
\raggedright
  \item \emph{Say it:} \emph{gambhīr}
  \item \emph{Say it:} \emph{gambhīr}, once more
  \item \emph{Say it:} say \emph{gambhīr}
  \item \emph{Recall:} say the Marathi for about, then the Marathi for only, then say \emph{gambhīr} again
\end{itemize}

\begin{tcolorbox}[breakable,colback=teal!4,colframe=teal!35!black,title={Before you move on}]
What does \mr{गंभीर} mean? (\textquotedblleft{}Serious\textquotedblright{}.) Say it once more.
\end{tcolorbox}
\section[\texorpdfstring{\mr{लाजाळू} (lājāḷū)}{लाजाळू (lājāḷū)}]{\mr{लाजाळू} (lājāḷū) — shy}

\label{lesson:MR-C204-lajalu}



\begin{quote}
Before the new one: say the Marathi for only, then the Marathi for serious.
\end{quote}

\subsection*{You'll want to know: \mr{लाजाळू}}
\textbf{\mr{लाजाळू}} — \emph{lājāḷū} — \textquotedblleft{}shy\textquotedblright{}.

\begin{cousinweb}
\textbf{\mr{लाजाळू}} (\emph{lājāḷū}) is \textbf{\mr{लाज}}, \textquotedblleft{}shame\textquotedblright{}, from Sanskrit \emph{lajjā}, \textquotedblleft{}shame, modesty\textquotedblright{} (CDIAL 10910), with the ending \emph{-āḷū}, \textquotedblleft{}given to\textquotedblright{}.
\end{cousinweb}

\begin{grammarlens}[title={What you've built}]
One more word for this chapter.
\end{grammarlens}

\subsection*{Your turn}
\begin{itemize}
\raggedright
  \item \emph{Say it:} \emph{lājāḷū}
  \item \emph{Say it:} \emph{lājāḷū}, once more
  \item \emph{Say it:} say \emph{lājāḷū}
  \item \emph{Recall:} say the Marathi for only, then the Marathi for serious, then say \emph{lājāḷū} again
\end{itemize}

\begin{tcolorbox}[breakable,colback=teal!4,colframe=teal!35!black,title={Before you move on}]
What does \mr{लाजाळू} mean? (\textquotedblleft{}Shy\textquotedblright{}.) Say it once more.
\end{tcolorbox}
\section[\texorpdfstring{\mr{शूर} (śūr)}{शूर (śūr)}]{\mr{शूर} (śūr) — valiant, brave}

\label{lesson:MR-C204-sur}



\begin{quote}
Before the new one: say the Marathi for serious, then the Marathi for shy.
\end{quote}

\subsection*{You'll want to know: \mr{शूर}}
\textbf{\mr{शूर}} — \emph{śūr} — \textquotedblleft{}valiant, brave\textquotedblright{}.

\begin{cousinweb}
\textbf{\mr{शूर}} (\emph{śūr}) is Sanskrit \emph{śūra}, \textquotedblleft{}strong, valiant; hero\textquotedblright{}, a tatsama (CDIAL 12569).
\end{cousinweb}

\begin{grammarlens}[title={What you've built}]
One more word for this chapter.
\end{grammarlens}

\subsection*{Your turn}
\begin{itemize}
\raggedright
  \item \emph{Say it:} \emph{śūr}
  \item \emph{Say it:} \emph{śūr}, once more
  \item \emph{Say it:} say \emph{śūr}
  \item \emph{Recall:} say the Marathi for serious, then the Marathi for shy, then say \emph{śūr} again
\end{itemize}

\begin{tcolorbox}[breakable,colback=teal!4,colframe=teal!35!black,title={Before you move on}]
What does \mr{शूर} mean? (\textquotedblleft{}Valiant, brave\textquotedblright{}.) Say it once more.
\end{tcolorbox}
\section[\texorpdfstring{\mr{भित्रा} / \mr{भित्री} (bhitrā / bhitrī)}{भित्रा / भित्री (bhitrā / bhitrī)}]{\mr{भित्रा} / \mr{भित्री} (bhitrā / bhitrī) — fearful, cowardly}

\label{lesson:MR-C204-bhitra}



\begin{quote}
Before the new one: say the Marathi for shy, then the Marathi for valiant.
\end{quote}

\subsection*{You'll want to know: \mr{भित्रा} / \mr{भित्री}}
\textbf{\mr{भित्रा} / \mr{भित्री}} — \emph{bhitrā / bhitrī} — \textquotedblleft{}fearful, cowardly\textquotedblright{}.

\begin{cousinweb}
The origin of \textbf{\mr{भित्रा}} (\emph{bhitrā}) is not settled. CDIAL has no entry for it; under Sanskrit \emph{bhīta}, \textquotedblleft{}frightened\textquotedblright{}, it lists only Old Marathi \emph{bheḍa}, \textquotedblleft{}timid\textquotedblright{} (CDIAL 9511).
\end{cousinweb}

\begin{grammarlens}[title={What you've built}]
One more word for this chapter.
\end{grammarlens}

\subsection*{Your turn}
\begin{itemize}
\raggedright
  \item \emph{Say it:} \emph{bhitrā / bhitrī}
  \item \emph{Say it:} \emph{bhitrā / bhitrī}, once more
  \item \emph{Say it:} say \emph{bhitrā / bhitrī}
  \item \emph{Recall:} say the Marathi for shy, then the Marathi for valiant, then say \emph{bhitrā / bhitrī} again
\end{itemize}

\begin{tcolorbox}[breakable,colback=teal!4,colframe=teal!35!black,title={Before you move on}]
What does \mr{भित्रा} / \mr{भित्री} mean? (\textquotedblleft{}Fearful, cowardly\textquotedblright{}.) Say it once more.
\end{tcolorbox}
\section[\texorpdfstring{\mr{उत्सुक} (utsuk)}{उत्सुक (utsuk)}]{\mr{उत्सुक} (utsuk) — curious, eager}

\label{lesson:MR-C204-utsuk}



\begin{quote}
Before the new one: say the Marathi for valiant, then the Marathi for fearful.
\end{quote}

\subsection*{You'll want to know: \mr{उत्सुक}}
\textbf{\mr{उत्सुक}} — \emph{utsuk} — \textquotedblleft{}curious, eager\textquotedblright{}.

\begin{cousinweb}
\textbf{\mr{उत्सुक}} (\emph{utsuk}) is Sanskrit \emph{utsuka}, \textquotedblleft{}eager, anxious\textquotedblright{}, a tatsama. CDIAL notes it, with a \textquotedblleft{}Cf.\textquotedblright{}, under a reconstructed \emph{*utsukya}, \textquotedblleft{}eagerness\textquotedblright{} (CDIAL 1886).
\end{cousinweb}

\begin{grammarlens}[title={What you've built}]
One more word for this chapter.
\end{grammarlens}

\subsection*{Your turn}
\begin{itemize}
\raggedright
  \item \emph{Say it:} \emph{utsuk}
  \item \emph{Say it:} \emph{utsuk}, once more
  \item \emph{Say it:} say \emph{utsuk}
  \item \emph{Recall:} say the Marathi for valiant, then the Marathi for fearful, then say \emph{utsuk} again
\end{itemize}

\begin{tcolorbox}[breakable,colback=teal!4,colframe=teal!35!black,title={Before you move on}]
What does \mr{उत्सुक} mean? (\textquotedblleft{}Curious, eager\textquotedblright{}.) Say it once more.
\end{tcolorbox}
